\documentclass[12pt, a4paper]{ctexart}
\usepackage{fontspec}
\usepackage{xcolor}
\usepackage{hyperref}
\usepackage{geometry}
\usepackage{enumitem}
\usepackage{parskip}
\usepackage{titlesec}
\usepackage{tcolorbox}
\tcbuselibrary{breakable}
\usepackage{setspace}
\usepackage{listings}

\geometry{left=2.5cm, right=2.5cm, top=2.5cm, bottom=2.5cm}

\setmainfont{Times New Roman}
\setCJKmainfont{SimSun}

\hypersetup{
    colorlinks=true,
    linkcolor=blue!70!black,
    urlcolor=cyan,
    pdftitle={Java泛型-逐题精讲解义},
    pdfauthor={专升本-fifine老师}
}

\definecolor{myblue}{RGB}{30,100,200}
\definecolor{lightblue}{RGB}{230,240,255}
\definecolor{darkblue}{RGB}{0,50,120}

\newtcolorbox{bluebox}[1][]{
    colback=lightblue,
    colframe=myblue,
    coltitle=darkblue,
    fonttitle=\bfseries,
    title={#1},
    boxrule=0.8pt,
    arc=4pt,
    left=6pt,
    right=6pt,
    top=6pt,
    bottom=6pt,
    before skip=8pt,
    after skip=8pt,
    breakable
}

\usepackage{etoolbox}
\BeforeBeginEnvironment{bluebox}{\par\vfill}%
\AfterEndEnvironment{bluebox}{\par\vfill}%

\titleformat{\section}
  {\normalfont\Large\bfseries\color{myblue}}{\thesection}{1em}{}
\titlespacing*{\section}{0pt}{3.5ex plus 1ex minus .2ex}{1.5ex plus .2ex}

\title{\textcolor{myblue}{\textbf{Java泛型 - 逐题精讲解义}}}
\author{专升本-fifine老师 教学组}
\date{2026年03月02日}

\lstset{
    language=Java,
    basicstyle=\ttfamily\small,
    keywordstyle=\color{blue},
    commentstyle=\color{green!60!black},
    stringstyle=\color{orange},
    numbers=left,
    numberstyle=\tiny\color{gray},
    frame=single,
    breaklines=true,
    columns=fullflexible,
    showstringspaces=false
}

\newcommand{\code}[1]{\texttt{#1}}

\begin{document}

\maketitle
\thispagestyle{empty}
\newpage

\section{泛型}

\begin{enumerate}[label=\textbf{\arabic*.}]

	\item \textbf{以下关于Java中的泛型的说法错误的是}
	      \begin{enumerate}[label=\Alph*.]
		      \item 可以提高代码的安全性
		      \item 可以减少类型转换
		      \item 泛型类型在运行时会被擦除
		      \item 泛型可以用于基本数据类型
	      \end{enumerate}

	      \begin{bluebox}{答案与深度解析}
		      \textbf{答案：D. 泛型可以用于基本数据类型}

		      \textbf{【核心认知框架】}

		      \textbf{【选项原理深度拆解】}
					  \begin{itemize}[label={}, leftmargin=*, nosep]


		      \item \textbf{A. 提高代码安全性 — 正确}
					  \begin{itemize}[label={}, nosep]
		      
			      \item 编译期类型检查
			      \item 减少ClassCastException
		      

					  \end{itemize}
		      \item \textbf{B. 减少类型转换 — 正确}
					  \begin{itemize}[label={}, nosep]
		      
			      \item 泛型不需要强制类型转换
			      \item List<String>直接get()就是String
		      

					  \end{itemize}
		      \item \textbf{C. 运行时类型擦除 — 正确}
					  \begin{itemize}[label={}, nosep]
		      
			      \item 泛型信息在编译后被擦除
			      \item 只保留原始类型(raw type)
			      \item 这是泛型的核心机制
		      

					  \end{itemize}
		      \item \textbf{D. 可用于基本数据类型 — 错误}
					  \begin{itemize}[label={}, nosep]
		      
			      \item 泛型不支持基本类型！
			      \item 必须使用包装类
			      \item 例如：不能写List<int>，只能写List<Integer>
		      

		      
					  \end{itemize}
					  \end{itemize}
\textbf{【应背诵知识点】}

		      泛型核心概念：
		      \begin{itemize}[label={}, leftmargin=*, nosep]
			      \item 类型擦除(Type Erasure)：编译后泛型信息被擦除为原始类型
			      \item 泛型限制：不能用于基本类型，只能用包装类
			      \item 泛型边界：<? extends T>、<? super T>、<?>
			      \item 泛型擦除后的类型：List<String> → List、List<Integer> → List
		      \end{itemize}

		      \textbf{【类型擦除详解】}

		      \begin{enumerate}[label={}]
			      \item 无边界泛型：Object
			            \begin{itemize}[label={}, leftmargin=*, nosep]
				            \item List<T> → List，元素类型变为Object
			            \end{itemize}
			      \item 有边界泛型：具体类型
			            \begin{itemize}[label={}, leftmargin=*, nosep]
				            \item List<? extends Number> → List（边界信息保存在Signature属性中）
			            \end{itemize}
			      \item 桥接方法：编译器自动生成
			            \begin{itemize}[label={}, leftmargin=*, nosep]
				            \item 用于保持多态兼容性
			            \end{itemize}
		      \end{enumerate}

		      \textbf{【命题陷阱破解】}

		      \begin{itemize}[label={}, leftmargin=*, nosep]
			      \item 陷阱1："泛型在运行时可用" → 错！运行时已擦除
			      \item 陷阱2："List<int>是合法的" → 错！必须用List<Integer>
			      \item 陷阱3："泛型ArrayList<String>和ArrayList<Object>可以互相转换" → 错！泛型类型不兼容
		      \end{itemize}

		      \textbf{【泛型 vs 数组】}

		      \begin{tabular}{|c|c|c|}
			      \hline
			      特性    & 泛型    & 数组        \\
			      \hline
			      类型安全  & 编译期检查 & 运行时检查     \\
			      \hline
			      元素类型  & 引用类型  & 基本类型+引用类型 \\
			      \hline
			      运行时信息 & 擦除    & 保留        \\
			      \hline
			      协变性   & 否     & 是         \\
			      \hline
		      \end{tabular}

		      \textbf{【实战代码示例】}

		      \begin{lstlisting}[language=Java]
// 正确：使用包装类
List<Integer> intList = new ArrayList<>();
intList.add(10); // 自动装箱
int num = intList.get(0); // 自动拆箱

// 错误：不能使用基本类型
// List<int> intList = new ArrayList<>(); // 编译错误！

// 类型擦除示例
List<String> stringList = new ArrayList<>();
List<Integer> intList2 = new ArrayList<>();
// 运行时：stringList和intList2的类对象相同
// 都是ArrayList
System.out.println(stringList.getClass() == intList2.getClass()); // true
\end{lstlisting}
	      \end{bluebox}

	      \vspace{1em}

	\item \textbf{以下哪个是正确的Java泛型声明？}
	      \begin{enumerate}[label=\Alph*.]
		      \item class MyClass<T extends Number>
		      \item class MyClass<T super Number>
		      \item class MyClass<T implements Number>
		      \item class MyClass<T>
	      \end{enumerate}

	      \begin{bluebox}{答案与深度解析}
		      \textbf{答案：A. class MyClass<T extends Number>}

		      \textbf{【核心认知框架】}

		      \textbf{【选项原理深度拆解】}
					  \begin{itemize}[label={}, leftmargin=*, nosep]


		      \item \textbf{A. <T extends Number> — 正确}
					  \begin{itemize}[label={}, nosep]
		      
			      \item 使用extends指定泛型上界
			      \item 表示T必须是Number或其子类
		      

					  \end{itemize}
		      \item \textbf{B. <T super Number> — 错误}
					  \begin{itemize}[label={}, nosep]
		      
			      \item super不能用于泛型类声明
			      \item super只用于泛型通配符：<? super Number>
		      

					  \end{itemize}
		      \item \textbf{C. <T implements Number> — 错误}
					  \begin{itemize}[label={}, nosep]
		      
			      \item implements不能用于泛型
			      \item 泛型上界使用extends
			      \item 注意：Number是类，不是接口
		      

					  \end{itemize}
		      \item \textbf{D. <T> — 正确但不完整}
					  \begin{itemize}[label={}, nosep]
		      
			      \item 语法正确，但题目要求"正确的泛型声明"
			      \item A选项展示了泛型边界的正确用法
		      

		      
					  \end{itemize}
					  \end{itemize}
\textbf{【应背诵知识点】}

		      泛型边界：
		      \begin{itemize}[label={}, leftmargin=*, nosep]
			      \item extends：指定上界（Producer Extends）
			            \begin{itemize}[label={}, leftmargin=*, nosep]
				            \item List<? extends Number>：可以是List<Integer>、List<Double>
				            \item 只能读取（作为生产者）
			            \end{itemize}
			      \item super：指定下界（Consumer Super）
			            \begin{itemize}[label={}, leftmargin=*, nosep]
				            \item List<? super Integer>：可以是List<Number>、List<Object>
				            \item 只能写入（作为消费者）
			            \end{itemize}
			      \item PECS原则：Producer Extends, Consumer Super
		      \end{itemize}

		      \textbf{【泛型边界 vs implements】}

		      \begin{enumerate}[label={}]
			      \item extends用于泛型边界
			            \begin{itemize}[label={}, leftmargin=*, nosep]
				            \item class MyClass<T extends Number> ✓
				            \item interface MyInterface<T extends Serializable> ✓
			            \end{itemize}
			      \item implements用于类实现接口
			            \begin{itemize}[label={}, leftmargin=*, nosep]
				            \item class MyClass implements MyInterface ✗（泛型中不能这样用）
			            \end{itemize}
			      \item 注意：即使Number是类，也用extends！
			            \begin{itemize}[label={}, leftmargin=*, nosep]
				            \item extends可以用于类和接口
			            \end{itemize}
		      \end{enumerate}

		      \textbf{【命题陷阱破解】}

		      \begin{itemize}[label={}, leftmargin=*, nosep]
			      \item 陷阱1："implements用于泛型" → 错！泛型边界用extends
			      \item 陷阱2："super用于泛型类声明" → 错！super只用于通配符
			      \item 陷阱3："Number是接口所以用implements" → 错！Number是类，也用extends
			      \item 陷阱4："没有边界声明就是错误的" → 错！<T>本身是合法的
		      \end{itemize}

		      \textbf{【通配符详解】}

		      \begin{enumerate}[label={}]
			      \item 无界通配符：<?>
			            \begin{itemize}[label={}, leftmargin=*, nosep]
				            \item List<?>：未知类型的列表
				            \item 可以读取Object，但不能写入
			            \end{itemize}
			      \item 上界通配符：<? extends T>
			            \begin{itemize}[label={}, leftmargin=*, nosep]
				            \item 适用于"读取"为主的场景
				            \item PECS中的"Producer"
			            \end{itemize}
			      \item 下界通配符：<? super T>
			            \begin{itemize}[label={}, leftmargin=*, nosep]
				            \item 适用于"写入"为主的场景
				            \item PECS中的"Consumer"
			            \end{itemize}
		      \end{enumerate}

		      \textbf{【实战代码示例】}

		      \begin{lstlisting}[language=Java]
// 正确：使用extends指定上界
class Container<T extends Number> {
    private T value;
    public void setValue(T value) { this.value = value; }
    public T getValue() { return value; }
}

// 使用
Container<Integer> c1 = new Container<>(); // OK
Container<Double> c2 = new Container<>();  // OK
// Container<String> c3 = new Container<>(); // 编译错误！String不是Number子类

// 正确：使用super指定下界
void addNumbers(List<? super Integer> list) {
    list.add(10); // 可以写入Integer
    // Integer num = list.get(0); // 错误！只能读取Object
}

// 错误：super不能用于类声明
// class MyClass<T super Number> {} // 编译错误！

// 错误：implements不能用于泛型
// class MyClass<T implements Number> {} // 编译错误！
\end{lstlisting}
	      \end{bluebox}

\end{enumerate}

\end{document}
